% v6 section. Every number is a generated \vsix macro (scripts/hero_v6.py,
% scripts/v6_macros.py); none is typed here. To be \input during the post-PR-6 integration pass.
\section{Stopping a Model Is Not Stopping Its Effects}\label{sec:v6stop}

An operator who decides to stop an agent can act at three places. \emph{Halting inference} ends the stream of new
proposals. \emph{Revoking admission} at the broker refuses permits for proposals that arrive later. \emph{Fencing the
sink} raises the epoch floor at the component that writes durable state, so that any operation carrying an older epoch
is refused at commit, including operations that were admitted, dispatched and are already running. A system that
tracks and stops model workloads~\cite{controlroom} acts at the first, on the producer of proposals. Only the third reaches work that is already in flight, and it
is the fencing-token argument of distributed locking~\cite{kleppmann2016} applied to agent authority.

\para{Setup (E4).} We run the real broker and sink processes (Rust, SQLite, JSON-lines RPC; Section~\ref{sec:core})
under a synthetic agent that emits \(\lambda = 2\) tool calls per second to \(C \in \{4, 32\}\) concurrent executors
whose latencies are exponential with mean \(w \in \{0.5, 5, 30\}\)\,s. Each call is a pre-approved, authorized one-unit
effect; nothing in the workload is adversarial. After a 20\,s warm-up the operator decides to stop, and the stop takes
effect \(L \in \{0, 1, 5, 35\}\)\,s later. A trial counts the effects committed after the stop took effect, timed
by the harness's log of the sink's replies (the set of committed effects must equal the sink's database), and the
legitimate calls the stop discards (cancelled or denied). A reply follows its commit, so this timing can place a commit
later than it happened but never earlier: it can add effects after the stop, never hide one. E4b, which kills the sink
inside requests, counts from the database instead. The plan, grid and
metrics were committed before the first trial; the \vsixEFourTrials{}
trials include \vsixEFourFaultTrials{} with SIGKILL of the sink or delayed completions.

\para{Result.} We draw the stable cell with the longest latency (\(C = \vsixStopCalls\),
\(w = \vsixStopLatency\)\,s, \(L = 0\)), chosen by that rule rather than by eye; the rule was written after E4's
data, not in its plan, and the grid-wide ranges follow. In E4, after the stop took effect, halting
inference let \vsixHaltRange{} effects commit and revoking at the broker \vsixRevokeRange{}, the last
\vsixLastCommitSeconds\,s later; fencing the sink let \vsixFenceRange{} commit and discarded \vsixFenceLost{}
legitimate calls. Across all \vsixEFourCells{} fault-free cells, halting leaked up to \vsixGridHaltMax{} effects and
revoking up to \vsixGridRevokeMax; across \vsixFenceTrials{} fence trials, the most effects any one committed after
the fence took effect was \vsixFencePostMax. \vsixEFourIdentitySentence

\para{Where the stop lands (E4c, post hoc).} E4's warm-up was a whole number of the agent's 2\,s emission cycles, so
the decision always fell just after a burst of four calls, which revoking at once refused and halting did not. After an
internal review pointed this out, we fixed a plan for one more run of the drawn cell with the decision at a uniformly
random point of the cycle, \vsixPhaseReplicates{} replicates per stop, the three stops of a replicate paired by seed
and phase, and committed the front figure to its result before running it. Halting let a median of
\vsixPhaseHaltMedian{} effects commit after the stop took effect (range \vsixPhaseHaltRange), revoking
\vsixPhaseRevokeMedian{} (\vsixPhaseRevokeRange) and fencing \vsixPhaseFenceRange, which discarded a median of
\vsixPhaseFenceLostMedian{} legitimate calls (\vsixPhaseFenceLostRange). Paired by replicate, halting and revoking
differed by \vsixPhaseDiffMedian{} in the median (range \vsixPhaseDiffRange; no difference in \vsixPhaseDiffZero{} of
\vsixPhaseReplicates). Figure~\ref{fig:v6hero}a draws this run. E4's gap between halting and revoking came from where
its stop landed: when calls are dispatched as soon as they are emitted, revoking at the broker refuses almost nothing
that halting does not, and both leave open the work already in flight.

\para{What the numbers mean.} The fence's zero is a property of the design, not a discovery: the sink serializes its
RPCs, \texttt{fence} cancels every pending record below the new floor, and effect identifiers are unique in the schema,
so duplicates are impossible rather than rare. What E4 measures is (i) that this implementation keeps the property
under the faults we injected, and (ii) the size and cost of the window the other two layers leave open. That window is
queueing, and it behaves like queueing. Where the offered load \(\rho = \lambda w / C\) is below one, halting leaks
the calls in flight and those emitted before the stop takes effect, which the plan predicted from Little's law.
\vsixLittleSentence{} Revocation leaks the calls already accepted.
Where \(\rho \ge 1\) (\vsixEFourOverloadedCells{} of the \vsixEFourWCCells{} \((w, C)\) cells), the backlog grows during the warm-up and
halting leaks all of it; we report those cells separately and do not apply Little's law to them. The fence pays for its
zero in the opposite currency. Every legitimate call in flight at the stop is discarded, which is what an operator who
stops an agent usually wants and what an operator who meant ``pause'' does not.

\para{Which term each stop removes.} Counting from the decision rather than from enforcement separates two terms
(Figure~\ref{fig:v6hero}b). During the delay \(L\) no layer acts, and about \(\lambda L\) effects commit at call rate
\(\lambda\). When the stop takes effect, the calls already dispatched add about \(\lambda w\) for halting and fewer for
revocation, which refuses those not yet accepted. The fence removes the second term and none of the first. In the drawn
cell at \(L = \vsixDelayMax\)\,s, with a measured call rate of \vsixDelayLambda{} per second, halting let
\vsixDelayHaltAtMax{} effects commit after the decision and fencing \vsixDelayFenceAtMax. Whenever enforcement is slow
compared with tool latency, shortening \(L\) is the larger lever; the fence is what makes the remainder zero. The plan
predicted \(\lambda(L+w)\) for halting; reading the fence's count as \(\lambda L\) is post hoc.

\para{Why not kill the executors?} An operator who controls every executor can stop each one before it commits, and
in this harness that would also leave zero, for the same reason the fence does: the executors are our processes and
every effect lands at our sink. We did not run that arm, because its zero would hold by construction as well. The fence
needs no such control. It acts at the one place every effect must pass, so it also covers an executor the operator
cannot reach (a remote runner, a third-party service, a host on the other side of a partition), a commit request
already sent when the kill arrives, and a retry after the sink restarts; E4b below measures the last two. Killing what
can be reached and fencing what must not commit are complementary, and the fence is the one that does not depend on
knowing where every executor is. The other natural strengthening, revoking at the broker and also invalidating the
permits it has already issued, is a fence exactly when something checks the invalidation at commit; checked anywhere
earlier, a permit already presented is not reached.

\para{Faults (E4b).} A SIGKILL that lands between requests tests little. In a separately predeclared run we kill the sink while a \texttt{complete} or a \texttt{fence} request is
outstanding, hold the executor's \texttt{accept} across the fence, and drive the system with closed-loop agents that
wait for each result. \vsixEFourBSentence

\para{Scope.} The sink mediates only its own SQLite effects. An external effect (a payment API, a repository, a cloud
resource) is fenced only if that surface can refuse a stale epoch, which is the question of
Section~\ref{sec:v6boundary} and Figure~\ref{fig:v6hero}c. No arm aborts an executor that is already running: the
fence refuses its commit, but anything an executor does outside the sink before it commits is stopped by none of the
three layers. The latencies are sampled, not production traces, and the agent is synthetic.
